\begin{PSbox}[unbreakable]{obj}{Learning Objectives}

After completing this module, students will be able to:

\begin{enumerate}
\def\labelenumi{\arabic{enumi}.}
\tightlist
\item
  Define random variables and explain their role in engineering analysis
\item
  Distinguish between discrete and continuous random variables
\item
  Compute and interpret PMF, PDF, and CDF
\item
  Clearly explain the difference between probability mass, density, and cumulative probability
\item
  Apply transformations to random variables
\item
  Use MATLAB to generate, visualize, and analyze random variables
\end{enumerate}

\end{PSbox}

\PSsection{2.1}{Introduction: Why Random Variables?}

\PSsubsection{The Problem}

In Module 1, we worked with events — subsets of a sample space. But engineering analysis requires \textbf{numbers}:

\begin{itemize}
\tightlist
\item
  How much noise voltage is present?
\item
  How many packets arrived?
\item
  What is the received signal power?
\end{itemize}

A \textbf{random variable} maps outcomes of a random experiment to real numbers, enabling us to use all the tools of calculus and algebra.

\PSsubsection{Engineering Motivation}

{\def\LTcaptype{none} % do not increment counter
\begin{longtable}[c]{@{}cc@{}}
\toprule\noalign{}
\rowcolor{PSnavy}\PSth{We want to analyze…} & \PSth{We define a random variable…} \\
\midrule\noalign{}
\endhead
\bottomrule\noalign{}
\endlastfoot
Noise in a circuit & $X$ = noise voltage (volts) \\
Bit errors in a frame & $N$ = number of bit errors \\
Component lifetime & $T$ = time until failure (hours) \\
Signal amplitude & A = received amplitude $(V)$ \\
Packet delay & $D$ = end-to-end delay $(ms)$ \\
\end{longtable}
}

\PSsection{2.2}{Definition of a Random Variable}

\PSsubsection{Formal Definition}

A \textbf{random variable} $X$ is a function that maps each outcome $\omega$ in the sample space $S$ to a real number:

\PSdisplay{X: S \rightarrow \mathbb{R}}\PSdisplay{\omega \mapsto X(\omega)}

\PSsubsection{Interpretation}

The random variable assigns a numerical value to each possible outcome of the experiment. The randomness comes from which outcome $\omega$ occurs — once $\omega$ is determined, $X(\omega)$ is a definite number.

\PSsubsection{Engineering Example}

\textbf{Experiment:} Transmit 5 bits over a noisy channel.\PSbr{}\textbf{Sample space:} $S$ = all possible sequences of correct/error for 5 bits ($2^{5} = 32$ outcomes)\PSbr{}\textbf{Random variable:} $X$ = number of bit errors

{\def\LTcaptype{none} % do not increment counter
\begin{longtable}[c]{@{}cc@{}}
\toprule\noalign{}
\rowcolor{PSnavy}\PSth{Outcome $\omega$} & \PSth{$X(\omega)$} \\
\midrule\noalign{}
\endhead
\bottomrule\noalign{}
\endlastfoot
$(C,\,C,\,C,\,C,\,C)$ & 0 \\
$(E,\,C,\,C,\,C,\,C)$ & 1 \\
$(E,\,E,\,C,\,C,\,C)$ & 2 \\
… & … \\
$(E,\,E,\,E,\,E,\,E)$ & 5 \\
\end{longtable}
}

Now instead of working with 32 outcomes, we work with a single number $X \in \{0,\,1,\,2,\,3,\,4,\,5\}$.

\PSsection{2.3}{Discrete Random Variables}

\begin{PSbox}[unbreakable]{def}{Definition}

A random variable $X$ is \textbf{discrete} if it takes values from a \textbf{countable} set (finite or countably infinite).

\end{PSbox}

\PSsubsection{Characteristics}

\begin{itemize}
\tightlist
\item
  Values can be listed: $\{x_{1},\, x_{2},\, x_{3},\, \ldots \}$
\item
  Each value has a specific probability (a ``mass'' of probability)
\item
  Sum of all probabilities = 1
\end{itemize}

\PSsubsection{Engineering Examples}

{\def\LTcaptype{none} % do not increment counter
\begin{longtable}[c]{@{}
  >{\centering\arraybackslash}p{(0.965\linewidth - 4\tabcolsep) * \real{0.3810}}
  >{\centering\arraybackslash}p{(0.965\linewidth - 4\tabcolsep) * \real{0.4048}}
  >{\centering\arraybackslash}p{(0.965\linewidth - 4\tabcolsep) * \real{0.2143}}@{}}
\toprule\noalign{}
\rowcolor{PSnavy}\PSth{Random Variable} & \PSth{Possible Values} & \PSth{Context} \\
\midrule\noalign{}
\endhead
\bottomrule\noalign{}
\endlastfoot
Number of bit errors & $\{0,\, 1,\, 2,\, \ldots,\, n\}$ & Packet transmission \\
Number of retransmissions & \{0, 1, 2, …\} & ARQ protocol \\
Quantized voltage level & $\{0,\, 1,\, \ldots,\, 2^{N}-1\}$ & N-bit ADC output \\
Number of arrivals per second & \{0, 1, 2, …\} & Network traffic \\
Number of defective chips & $\{0,\, 1,\, \ldots,\, n\}$ & Wafer testing \\
\end{longtable}
}

\PSsection{2.4}{Continuous Random Variables}

\begin{PSbox}[unbreakable]{def}{Definition}

A random variable $X$ is \textbf{continuous} if it takes values from an \textbf{uncountable} set (interval of real numbers).

\end{PSbox}

\PSsubsection{Characteristics}

\begin{itemize}
\tightlist
\item
  $X$ can take any value in a range (e.g., $[0,\, \infty)$ or $(-\infty,\, \infty)$)
\item
  Probability of any single exact value is ZERO: $P(X = x) = 0$
\item
  Probabilities are defined for intervals: $P(a \le X \le b)$
\end{itemize}

\PSsubsection{Engineering Examples}

{\def\LTcaptype{none} % do not increment counter
\begin{longtable}[c]{@{}
  >{\centering\arraybackslash}p{(0.965\linewidth - 4\tabcolsep) * \real{0.5000}}
  >{\centering\arraybackslash}p{(0.965\linewidth - 4\tabcolsep) * \real{0.2188}}
  >{\centering\arraybackslash}p{(0.965\linewidth - 4\tabcolsep) * \real{0.2812}}@{}}
\toprule\noalign{}
\rowcolor{PSnavy}\PSth{Random Variable} & \PSth{Range} & \PSth{Context} \\
\midrule\noalign{}
\endhead
\bottomrule\noalign{}
\endlastfoot
Thermal noise voltage & $(-\infty,\, \infty)$ & Circuit noise \\
Time until failure & $[0,\, \infty)$ & Reliability \\
Phase of received signal & $[0,\, 2\pi)$ & Communications \\
Resistor value & $(R_{0} - \delta,\, R_{0} + \delta)$ & Manufacturing \\
Channel gain & $[0,\, \infty)$ & Wireless fading \\
\end{longtable}
}

\PSsubsection{Why $P(X = x) = 0$ for Continuous RVs}

Imagine measuring a voltage. The probability that $V$ = exactly 3.14159265… volts (infinite precision) is zero. We can only ask about intervals: $P(3.14 < V < 3.15)$.

\PSsection{2.5}{Probability Mass Function (PMF)}

\begin{PSbox}[unbreakable]{def}{Definition}

For a discrete random variable $X$, the \textbf{probability mass function} is:

\PSdisplay{p_X(x) = P(X = x)}

\end{PSbox}

\PSsubsection{Properties}

\begin{enumerate}
\def\labelenumi{\arabic{enumi}.}
\tightlist
\item
  \textbf{Non-negativity:} $p_{X}(x) \ge 0$ for all $x$
\item
  \textbf{Normalization:} $\sum p_{X}(x_{i}) = 1$ (sum over all possible values)
\item
  \textbf{Probability computation:} $P(X \in A) = \sum_{x \in A} p_{X}(x)$
\end{enumerate}

\PSsubsection{Visualization}

The PMF is represented as \textbf{vertical lines} (impulses) at each possible value, with height equal to the probability.

\PSfigure{m02-pmf-stems}

\PSsubsection{Engineering Example: Bit Errors in a Byte}

Transmit 8 bits, each with independent error probability $p = 0.1$.\PSbr{}$X$ = number of bit errors.

$p_{X}(k) = C(8,\,k) \cdot (0.1)^{k} \cdot (0.9)^{8-k}$, for $k = 0,\, 1,\, \ldots,\, 8$

{\def\LTcaptype{none} % do not increment counter
\begin{longtable}[c]{@{}cc@{}}
\toprule\noalign{}
\rowcolor{PSnavy}\PSth{$k$} & \PSth{$p_{X}(k)$} \\
\midrule\noalign{}
\endhead
\bottomrule\noalign{}
\endlastfoot
0 & 0.4305 \\
1 & 0.3826 \\
2 & 0.1488 \\
3 & 0.0331 \\
4 & 0.0046 \\
5+ & \textless{} 0.001 \\
\end{longtable}
}

\PSsection{2.6}{Probability Density Function (PDF)}

\begin{PSbox}[unbreakable]{def}{Definition}

For a continuous random variable $X$, the \textbf{probability density function} $f_{X}(x)$ is defined such that:

\PSdisplay{P(a \leq X \leq b) = \int_a^b f_X(x) \, dx}

\end{PSbox}

\PSsubsection{Properties}

\begin{enumerate}
\def\labelenumi{\arabic{enumi}.}
\tightlist
\item
  \textbf{Non-negativity:} $f_{X}(x) \ge 0$ for all $x$
\item
  \textbf{Normalization:} $\int_{-\infty}^{\infty} f_{X}(x)\,dx = 1$
\item
  \textbf{Probability = Area under curve} between two points
\end{enumerate}

\PSsubsection{Critical Point: $f_{X}(x)$ is NOT a Probability!}

\begin{PSquote}

\PSwarn{} $f_{X}(x)$ can be GREATER than 1!

\end{PSquote}

The PDF gives probability \textbf{density} — probability per unit length. Only the \textbf{integral} (area) gives probability.

\textbf{Analogy:} Mass density (kg/m\textsuperscript{3}) is not mass. You need to integrate density over volume to get mass. Similarly, probability density is not probability — you integrate over an interval to get probability.

\PSsubsection{Visualization}

\PSfigure{m02-pdf-area}

\begin{PSbox}[unbreakable]{ex}{Engineering Example: Thermal Noise Voltage}

Thermal noise voltage in a circuit often follows a Gaussian distribution:

\PSdisplay{f_X(x) = \frac{1}{\sigma\sqrt{2\pi}} e^{-x^2/(2\sigma^2)}}

where $\sigma$ is the RMS noise voltage.

$P(\text{noise exceeds 2}\sigma) = P(\lvert X\rvert > 2\sigma) = 1 - P(-2\sigma \le X \le 2\sigma) \approx 1 - 0.9545 = 0.0455$

\end{PSbox}

\PSsection{2.7}{Cumulative Distribution Function (CDF)}

\begin{PSbox}[unbreakable]{def}{Definition}

For any random variable $X$ (discrete or continuous), the \textbf{cumulative distribution function} is:

\PSdisplay{F_X(x) = P(X \leq x)}

\end{PSbox}

\PSsubsection{Properties}

\begin{enumerate}
\def\labelenumi{\arabic{enumi}.}
\tightlist
\item
  \textbf{Non-decreasing:} if $a < b$, then $F_{X}(a) \le F_{X}(b)$
\item
  \textbf{Limits:} $F_{X}(-\infty) = 0,\, F_{X}(+\infty) = 1$
\item
  \textbf{Right-continuous:} $F_{X}(x)$ is continuous from the right
\item
  \textbf{Probability of interval:} $P(a < X \le b) = F_{X}(b) - F_{X}(a)$
\end{enumerate}

\PSsubsection{CDF for Discrete RVs}

The CDF is a \textbf{staircase function} with jumps at each possible value:

\PSfigure{m02-cdf-discrete}

Jump at $x = k$ has height $p_{X}(k)$.

\PSsubsection{CDF for Continuous RVs}

The CDF is a \textbf{smooth, continuous} curve:

\PSdisplay{F_X(x) = \int_{-\infty}^{x} f_X(t) \, dt}

And conversely:\PSdisplay{f_X(x) = \frac{dF_X(x)}{dx}}

\PSfigure{m02-cdf-continuous}

\PSsubsection{Relationship: $\mathrm{PDF}$ \ensuremath{\leftrightarrow} CDF}

{\def\LTcaptype{none} % do not increment counter
\begin{longtable}[c]{@{}
  >{\centering\arraybackslash}p{(0.965\linewidth - 4\tabcolsep) * \real{0.3793}}
  >{\centering\arraybackslash}p{(0.965\linewidth - 4\tabcolsep) * \real{0.3103}}
  >{\centering\arraybackslash}p{(0.965\linewidth - 4\tabcolsep) * \real{0.3103}}@{}}
\toprule\noalign{}
\rowcolor{PSnavy}\PSth{Direction} & \PSth{Formula} & \PSth{Meaning} \\
\midrule\noalign{}
\endhead
\bottomrule\noalign{}
\endlastfoot
$\mathrm{PDF}$ \ensuremath{\to} CDF & $F_{X}(x) = \int_{-\infty}^{x} f_{X}(t)dt$ & Integrate density to get cumulative probability \\
$\mathrm{CDF}$ \ensuremath{\to} PDF & $f_{X}(x) = \frac{dF_{X}}{dx}$ & Differentiate CDF to get density \\
\end{longtable}
}

\PSsection{2.8}{KEY DISTINCTION: Mass vs. Density vs. Cumulative}

This is one of the most important conceptual distinctions in probability theory.

\PSsubsection{Comparison Table}

{\def\LTcaptype{none} % do not increment counter
\begin{longtable}[c]{@{}
  >{\centering\arraybackslash}p{(0.965\linewidth - 6\tabcolsep) * \real{0.2326}}
  >{\centering\arraybackslash}p{(0.965\linewidth - 6\tabcolsep) * \real{0.2558}}
  >{\centering\arraybackslash}p{(0.965\linewidth - 6\tabcolsep) * \real{0.2558}}
  >{\centering\arraybackslash}p{(0.965\linewidth - 6\tabcolsep) * \real{0.2558}}@{}}
\toprule\noalign{}
\rowcolor{PSnavy}\PSth{Property} & \PSth{PMF $p_{X}(x)$} & \PSth{PDF $f_{X}(x)$} & \PSth{CDF $F_{X}(x)$} \\
\midrule\noalign{}
\endhead
\bottomrule\noalign{}
\endlastfoot
\textbf{Applies to} & Discrete RVs only & Continuous RVs only & Both \\
\textbf{Gives directly} & $P(X = x)$ & Probability density at $x$ & $P(X \le x)$ \\
\textbf{Range of values} & [0, 1] & $[0,\, \infty)$ — can exceed 1! & [0, 1] \\
\textbf{Sum/Integral} & $\sum p_{X}(x) = 1$ & $\int f_{X}(x)dx = 1$ & Starts at 0, ends at 1 \\
\textbf{Get probability} & $P(X = x) = p_{X}(x)$ & $P(a \le X \le b) = \int f_{X}(x)dx$ & $P(a < X \le b) = F(b)-F(a)$ \\
\end{longtable}
}

\PSsubsection{Visual Analogy: The Warehouse}

Imagine probability as ``material'' distributed across possible values:

\textbf{Discrete (PMF):} Like storing material in numbered boxes. Each box has a specific mass. The PMF tells you the mass in each box.

\PSfigure{m02-pmf-boxes}

\textbf{Continuous (PDF):} Like spreading material along a continuous shelf. The density tells you how thick the material is at each point. You find mass by computing area (thickness \ensuremath{\times} width).

\PSfigure{m02-pdf-area-12}

\textbf{Cumulative (CDF):} Like a running total — how much material have you accumulated up to this point?

\PSsubsection{Why This Matters in Engineering}

\textbf{Example:} A noise voltage $X$ has PDF $f_{X}(x) = 2$ for $x \in [0,\, 0.5]$.

\begin{itemize}
\tightlist
\item
  $f_{X}(0.3) = 2$ — this is NOT a probability! It means the density is $2\,\mathrm{V}^{-1}$.
\item
  $P(0.2 \le X \le 0.4) = \int_{0.2}^{0.4} 2\,dx = 2 \times 0.2 = 0.4$ — THIS is a probability.
\item
  $P(X = 0.3) = 0$ — a single point has zero probability for continuous RVs.
\end{itemize}

\PSsubsection{Common Errors to Avoid}

{\def\LTcaptype{none} % do not increment counter
\begin{longtable}[c]{@{}
  >{\centering\arraybackslash}p{(0.965\linewidth - 2\tabcolsep) * \real{0.5000}}
  >{\centering\arraybackslash}p{(0.965\linewidth - 2\tabcolsep) * \real{0.5000}}@{}}
\toprule\noalign{}
\rowcolor{PSnavy}\PSth{\PSno{} Wrong Statement} & \PSth{\PSyes{} Correct Statement} \\
\midrule\noalign{}
\endhead
\bottomrule\noalign{}
\endlastfoot
``The probability at $x = 2$ is $f(2)$'' & ``The probability density at $x = 2$ is $f(2)$'' \\
``$f(x)$ must be \ensuremath{\le} 1'' & ``$f(x)$ can be any non-negative value'' \\
``$P(X = 3) = f(3)$ for continuous $X$'' & ``$P(X = 3) = 0$ for continuous $X$'' \\
``PMF and PDF are the same thing'' & ``PMF gives mass; PDF gives density'' \\
\end{longtable}
}

\PSsection{2.9}{Transformations of Random Variables}

\PSsubsection{Motivation}

In engineering, we often know the distribution of one quantity but need the distribution of a \textbf{function} of that quantity:

\begin{itemize}
\tightlist
\item
  Know voltage $V$ \ensuremath{\to} need power $P = \frac{V^{2}}{R}$
\item
  Know linear signal $X$ \ensuremath{\to} need dB value $Y = 10 \cdot \log_{10}(X)$
\item
  Know noise $N$ \ensuremath{\to} need $\lvert N\rvert$ (rectified noise)
\end{itemize}

\begin{PSbox}[unbreakable]{prob}{Problem Statement}

Given: $X$ with known PDF $f_{X}(x)$ or PMF $p_{X}(x)$\PSbr{}Find: The distribution of $Y = g(X)$

\end{PSbox}

\PSsubsection{Case 1: Discrete RV Transformation}

If $X$ is discrete with PMF $p_{X}(x)$, and $Y = g(X)$:

\PSdisplay{p_Y(y) = \sum_{x: g(x)=y} p_X(x)}

Sum the probabilities of all $x$ values that map to the same $y$.

\textbf{Example:} $X \in \{-2,\, -1,\, 0,\, 1,\, 2\}$ with equal probability 0.2 each.\PSbr{}$Y = X^{2}$

{\def\LTcaptype{none} % do not increment counter
\begin{longtable}[c]{@{}ccc@{}}
\toprule\noalign{}
\rowcolor{PSnavy}\PSth{$Y = X^{2}$} & \PSth{$X$ values mapping to $Y$} & \PSth{$p_{Y}(y)$} \\
\midrule\noalign{}
\endhead
\bottomrule\noalign{}
\endlastfoot
0 & \{0\} & 0.2 \\
1 & \{-1, 1\} & 0.4 \\
4 & \{-2, 2\} & 0.4 \\
\end{longtable}
}

\PSsubsection{Case 2: Monotonic Transformation (Continuous)}

If $Y = g(X)$ where $g$ is strictly monotonic and differentiable:

\PSdisplay{f_Y(y) = f_X(g^{-1}(y)) \cdot \left|\frac{dg^{-1}(y)}{dy}\right|}

Or equivalently:\PSdisplay{f_Y(y) = \frac{f_X(x)}{|g'(x)|}\bigg|_{x=g^{-1}(y)}}

\textbf{Derivation via CDF method:}

\begin{enumerate}
\def\labelenumi{\arabic{enumi}.}
\tightlist
\item
  Start with CDF: $F_{Y}(y) = P(Y \le y) = P(g(X) \le y)$
\item
  Invert: $= P(X \le g^{-1}(y))$ if $g$ is increasing
\item
  Differentiate: $f_{Y}(y) = f_{X}(g^{-1}(y)) \cdot \frac{d}{dy}[g^{-1}(y)]$
\end{enumerate}

\PSsubsection{Case 3: Non-Monotonic Transformation}

If $g$ is not monotonic, split the domain into intervals where $g$ IS monotonic:

\PSdisplay{f_Y(y) = \sum_i \frac{f_X(x_i)}{|g'(x_i)|}}

where $x_{1},\, x_{2}$, … are all solutions of $g(x) = y$.

\begin{PSbox}[unbreakable]{ex}{Engineering Example: Power from Voltage}

A noise voltage $X$ has PDF $f_{X}(x)$ (e.g., Gaussian with zero mean).\PSbr{}Power dissipated in a 1\ensuremath{\Omega} resistor: $P = X^{2}$

Since $Y = X^{2}$ is NOT monotonic (both $+x$ and $-x$ give the same $y$):

For $y > 0$: $x = +\sqrt{y}$ and $x = -\sqrt{y}$ are both solutions.\PSbr{}$g(x) = x^{2}$, so $g'(x) = 2x$

\PSdisplay{f_Y(y) = \frac{f_X(\sqrt{y})}{2\sqrt{y}} + \frac{f_X(-\sqrt{y})}{2\sqrt{y}}, \quad y > 0}

If $X$ is zero-mean Gaussian: $f_{X}(x) = f_{X}(-x)$, so:\PSdisplay{f_Y(y) = \frac{f_X(\sqrt{y})}{\sqrt{y}}, \quad y > 0}

\end{PSbox}

\begin{PSbox}[unbreakable]{ex}{Engineering Example: Linear Amplification}

Signal $Y$ = aX $+ b$ (amplifier with gain a and offset b).\PSbr{}$X$ has PDF $f_{X}(x)$.

$g^{-1}(y) = \frac{y-b}{a}$, and $\lvert \frac{d}{dy}\left[\frac{y-b}{a}\right]\rvert = \frac{1}{\lvert a\rvert}$

\PSdisplay{f_Y(y) = \frac{1}{|a|} f_X\left(\frac{y-b}{a}\right)}

\end{PSbox}

\PSsection{2.10}{MATLAB Examples}

\begin{PSbox}[unbreakable]{ex}{Example 1: Plotting a PMF}

\tcbset{PScodemode/.style={unbreakable}}

\begin{Shaded}
\begin{Highlighting}[]
\CommentTok{\%\% PMF of Number of Bit Errors in a Byte}
\VariableTok{n} \OperatorTok{=} \FloatTok{8}\OperatorTok{;}           \CommentTok{\% bits per byte}
\VariableTok{p} \OperatorTok{=} \FloatTok{0.1}\OperatorTok{;}         \CommentTok{\% bit error probability}
\VariableTok{k} \OperatorTok{=} \FloatTok{0}\OperatorTok{:}\VariableTok{n}\OperatorTok{;}         \CommentTok{\% possible values}

\CommentTok{\% Compute PMF (binomial distribution)}
\VariableTok{pmf} \OperatorTok{=} \VariableTok{binopdf}\NormalTok{(}\VariableTok{k}\OperatorTok{,} \VariableTok{n}\OperatorTok{,} \VariableTok{p}\NormalTok{)}\OperatorTok{;}

\CommentTok{\% Plot}
\VariableTok{figure}\OperatorTok{;}
\VariableTok{stem}\NormalTok{(}\VariableTok{k}\OperatorTok{,} \VariableTok{pmf}\OperatorTok{,} \SpecialStringTok{\textquotesingle{}filled\textquotesingle{}}\OperatorTok{,} \SpecialStringTok{\textquotesingle{}LineWidth\textquotesingle{}}\OperatorTok{,} \FloatTok{2}\NormalTok{)}\OperatorTok{;}
\VariableTok{xlabel}\NormalTok{(}\SpecialStringTok{\textquotesingle{}Number of Errors (k)\textquotesingle{}}\NormalTok{)}\OperatorTok{;}
\VariableTok{ylabel}\NormalTok{(}\SpecialStringTok{\textquotesingle{}P(X = k)\textquotesingle{}}\NormalTok{)}\OperatorTok{;}
\VariableTok{title}\NormalTok{(}\SpecialStringTok{\textquotesingle{}PMF: Bit Errors in 8{-}bit Byte (p = 0.1)\textquotesingle{}}\NormalTok{)}\OperatorTok{;}
\VariableTok{grid} \VariableTok{on}\OperatorTok{;}

\CommentTok{\% Verify normalization}
\VariableTok{fprintf}\NormalTok{(}\SpecialStringTok{\textquotesingle{}Sum of PMF = \%.6f (should be 1)\textbackslash{}n\textquotesingle{}}\OperatorTok{,} \VariableTok{sum}\NormalTok{(}\VariableTok{pmf}\NormalTok{))}\OperatorTok{;}
\end{Highlighting}
\end{Shaded}

\end{PSbox}

\begin{PSbox}[unbreakable]{ex}{Example 2: Plotting a PDF and Verifying Properties}

\tcbset{PScodemode/.style={unbreakable}}

\begin{Shaded}
\begin{Highlighting}[]
\CommentTok{\%\% PDF of Gaussian Noise Voltage}
\VariableTok{sigma} \OperatorTok{=} \FloatTok{0.5}\OperatorTok{;}     \CommentTok{\% RMS noise voltage}
\VariableTok{x} \OperatorTok{=} \VariableTok{linspace}\NormalTok{(}\OperatorTok{{-}}\FloatTok{3}\OperatorTok{*}\VariableTok{sigma}\OperatorTok{,} \FloatTok{3}\OperatorTok{*}\VariableTok{sigma}\OperatorTok{,} \FloatTok{1000}\NormalTok{)}\OperatorTok{;}

\CommentTok{\% Compute PDF}
\VariableTok{pdf\_x} \OperatorTok{=} \VariableTok{normpdf}\NormalTok{(}\VariableTok{x}\OperatorTok{,} \FloatTok{0}\OperatorTok{,} \VariableTok{sigma}\NormalTok{)}\OperatorTok{;}

\CommentTok{\% Plot}
\VariableTok{figure}\OperatorTok{;}
\VariableTok{plot}\NormalTok{(}\VariableTok{x}\OperatorTok{,} \VariableTok{pdf\_x}\OperatorTok{,} \SpecialStringTok{\textquotesingle{}b{-}\textquotesingle{}}\OperatorTok{,} \SpecialStringTok{\textquotesingle{}LineWidth\textquotesingle{}}\OperatorTok{,} \FloatTok{2}\NormalTok{)}\OperatorTok{;}
\VariableTok{xlabel}\NormalTok{(}\SpecialStringTok{\textquotesingle{}Voltage (V)\textquotesingle{}}\NormalTok{)}\OperatorTok{;}
\VariableTok{ylabel}\NormalTok{(}\SpecialStringTok{\textquotesingle{}f\_X(x) [V\^{}\{{-}1\}]\textquotesingle{}}\NormalTok{)}\OperatorTok{;}
\VariableTok{title}\NormalTok{(}\SpecialStringTok{\textquotesingle{}PDF: Gaussian Noise (σ = 0.5 V)\textquotesingle{}}\NormalTok{)}\OperatorTok{;}
\VariableTok{grid} \VariableTok{on}\OperatorTok{;}

\CommentTok{\% Verify normalization (numerical integration)}
\VariableTok{dx} \OperatorTok{=} \VariableTok{x}\NormalTok{(}\FloatTok{2}\NormalTok{) }\OperatorTok{{-}} \VariableTok{x}\NormalTok{(}\FloatTok{1}\NormalTok{)}\OperatorTok{;}
\VariableTok{total\_area} \OperatorTok{=} \VariableTok{trapz}\NormalTok{(}\VariableTok{x}\OperatorTok{,} \VariableTok{pdf\_x}\NormalTok{)}\OperatorTok{;}
\VariableTok{fprintf}\NormalTok{(}\SpecialStringTok{\textquotesingle{}Total area under PDF = \%.6f (should be 1)\textbackslash{}n\textquotesingle{}}\OperatorTok{,} \VariableTok{total\_area}\NormalTok{)}\OperatorTok{;}

\CommentTok{\% Note: PDF values can exceed 1!}
\VariableTok{fprintf}\NormalTok{(}\SpecialStringTok{\textquotesingle{}Maximum PDF value = \%.4f (can be \textgreater{} 1!)\textbackslash{}n\textquotesingle{}}\OperatorTok{,} \VariableTok{max}\NormalTok{(}\VariableTok{pdf\_x}\NormalTok{))}\OperatorTok{;}

\CommentTok{\% Compute P(|X| \textgreater{} 1V)}
\VariableTok{prob\_exceed} \OperatorTok{=} \FloatTok{1} \OperatorTok{{-}} \VariableTok{normcdf}\NormalTok{(}\FloatTok{1}\OperatorTok{,} \FloatTok{0}\OperatorTok{,} \VariableTok{sigma}\NormalTok{) }\OperatorTok{+} \VariableTok{normcdf}\NormalTok{(}\OperatorTok{{-}}\FloatTok{1}\OperatorTok{,} \FloatTok{0}\OperatorTok{,} \VariableTok{sigma}\NormalTok{)}\OperatorTok{;}
\VariableTok{fprintf}\NormalTok{(}\SpecialStringTok{\textquotesingle{}P(|X| \textgreater{} 1V) = \%.4f\textbackslash{}n\textquotesingle{}}\OperatorTok{,} \VariableTok{prob\_exceed}\NormalTok{)}\OperatorTok{;}
\end{Highlighting}
\end{Shaded}

\end{PSbox}

\begin{PSbox}[unbreakable]{ex}{Example 3: CDF Visualization}

\tcbset{PScodemode/.style={unbreakable}}

\begin{Shaded}
\begin{Highlighting}[]
\CommentTok{\%\% Comparing CDF for Discrete vs Continuous}

\VariableTok{figure}\OperatorTok{;}

\CommentTok{\% Subplot 1: Discrete CDF (Poisson, packets per second)}
\VariableTok{subplot}\NormalTok{(}\FloatTok{1}\OperatorTok{,}\FloatTok{2}\OperatorTok{,}\FloatTok{1}\NormalTok{)}\OperatorTok{;}
\VariableTok{lambda} \OperatorTok{=} \FloatTok{3}\OperatorTok{;}      \CommentTok{\% average arrivals per second}
\VariableTok{k} \OperatorTok{=} \FloatTok{0}\OperatorTok{:}\FloatTok{10}\OperatorTok{;}
\VariableTok{cdf\_discrete} \OperatorTok{=} \VariableTok{poisscdf}\NormalTok{(}\VariableTok{k}\OperatorTok{,} \VariableTok{lambda}\NormalTok{)}\OperatorTok{;}
\VariableTok{stairs}\NormalTok{(}\VariableTok{k}\OperatorTok{,} \VariableTok{cdf\_discrete}\OperatorTok{,} \SpecialStringTok{\textquotesingle{}r{-}\textquotesingle{}}\OperatorTok{,} \SpecialStringTok{\textquotesingle{}LineWidth\textquotesingle{}}\OperatorTok{,} \FloatTok{2}\NormalTok{)}\OperatorTok{;}
\VariableTok{xlabel}\NormalTok{(}\SpecialStringTok{\textquotesingle{}Number of Packets\textquotesingle{}}\NormalTok{)}\OperatorTok{;}
\VariableTok{ylabel}\NormalTok{(}\SpecialStringTok{\textquotesingle{}F\_X(x) = P(X ≤ x)\textquotesingle{}}\NormalTok{)}\OperatorTok{;}
\VariableTok{title}\NormalTok{(}\SpecialStringTok{\textquotesingle{}CDF: Discrete (Poisson, λ=3)\textquotesingle{}}\NormalTok{)}\OperatorTok{;}
\VariableTok{grid} \VariableTok{on}\OperatorTok{;} \VariableTok{ylim}\NormalTok{([}\FloatTok{0} \FloatTok{1.1}\NormalTok{])}\OperatorTok{;}

\CommentTok{\% Subplot 2: Continuous CDF (Exponential, time to failure)}
\VariableTok{subplot}\NormalTok{(}\FloatTok{1}\OperatorTok{,}\FloatTok{2}\OperatorTok{,}\FloatTok{2}\NormalTok{)}\OperatorTok{;}
\VariableTok{mu} \OperatorTok{=} \FloatTok{5}\OperatorTok{;}          \CommentTok{\% mean time to failure (hours)}
\VariableTok{t} \OperatorTok{=} \VariableTok{linspace}\NormalTok{(}\FloatTok{0}\OperatorTok{,} \FloatTok{20}\OperatorTok{,} \FloatTok{500}\NormalTok{)}\OperatorTok{;}
\VariableTok{cdf\_continuous} \OperatorTok{=} \VariableTok{expcdf}\NormalTok{(}\VariableTok{t}\OperatorTok{,} \VariableTok{mu}\NormalTok{)}\OperatorTok{;}
\VariableTok{plot}\NormalTok{(}\VariableTok{t}\OperatorTok{,} \VariableTok{cdf\_continuous}\OperatorTok{,} \SpecialStringTok{\textquotesingle{}b{-}\textquotesingle{}}\OperatorTok{,} \SpecialStringTok{\textquotesingle{}LineWidth\textquotesingle{}}\OperatorTok{,} \FloatTok{2}\NormalTok{)}\OperatorTok{;}
\VariableTok{xlabel}\NormalTok{(}\SpecialStringTok{\textquotesingle{}Time (hours)\textquotesingle{}}\NormalTok{)}\OperatorTok{;}
\VariableTok{ylabel}\NormalTok{(}\SpecialStringTok{\textquotesingle{}F\_X(x) = P(X ≤ x)\textquotesingle{}}\NormalTok{)}\OperatorTok{;}
\VariableTok{title}\NormalTok{(}\SpecialStringTok{\textquotesingle{}CDF: Continuous (Exponential, μ=5)\textquotesingle{}}\NormalTok{)}\OperatorTok{;}
\VariableTok{grid} \VariableTok{on}\OperatorTok{;} \VariableTok{ylim}\NormalTok{([}\FloatTok{0} \FloatTok{1.1}\NormalTok{])}\OperatorTok{;}
\end{Highlighting}
\end{Shaded}

\end{PSbox}

\PSsubsection{Example 4: Transformation of a Random Variable}

\tcbset{PScodemode/.style={unbreakable}}

\begin{Shaded}
\begin{Highlighting}[]
\CommentTok{\%\% Transformation: Power from Noise Voltage}
\CommentTok{\% X = noise voltage, Gaussian(0, sigma)}
\CommentTok{\% Y = X\^{}2 = instantaneous power (in 1{-}ohm load)}

\VariableTok{sigma} \OperatorTok{=} \FloatTok{1}\OperatorTok{;}
\VariableTok{N} \OperatorTok{=} \FloatTok{100000}\OperatorTok{;}

\CommentTok{\% Generate Gaussian noise samples}
\VariableTok{X} \OperatorTok{=} \VariableTok{sigma} \OperatorTok{*} \VariableTok{randn}\NormalTok{(}\FloatTok{1}\OperatorTok{,} \VariableTok{N}\NormalTok{)}\OperatorTok{;}

\CommentTok{\% Transform: power = voltage\^{}2}
\VariableTok{Y} \OperatorTok{=} \VariableTok{X}\OperatorTok{.\^{}}\FloatTok{2}\OperatorTok{;}

\CommentTok{\% Plot empirical PDF of Y using histogram}
\VariableTok{figure}\OperatorTok{;}
\VariableTok{subplot}\NormalTok{(}\FloatTok{2}\OperatorTok{,}\FloatTok{1}\OperatorTok{,}\FloatTok{1}\NormalTok{)}\OperatorTok{;}
\VariableTok{histogram}\NormalTok{(}\VariableTok{X}\OperatorTok{,} \FloatTok{100}\OperatorTok{,} \SpecialStringTok{\textquotesingle{}Normalization\textquotesingle{}}\OperatorTok{,} \SpecialStringTok{\textquotesingle{}pdf\textquotesingle{}}\NormalTok{)}\OperatorTok{;}
\VariableTok{hold} \VariableTok{on}\OperatorTok{;}
\VariableTok{x\_theory} \OperatorTok{=} \VariableTok{linspace}\NormalTok{(}\OperatorTok{{-}}\FloatTok{4}\OperatorTok{,} \FloatTok{4}\OperatorTok{,} \FloatTok{500}\NormalTok{)}\OperatorTok{;}
\VariableTok{plot}\NormalTok{(}\VariableTok{x\_theory}\OperatorTok{,} \VariableTok{normpdf}\NormalTok{(}\VariableTok{x\_theory}\OperatorTok{,} \FloatTok{0}\OperatorTok{,} \VariableTok{sigma}\NormalTok{)}\OperatorTok{,} \SpecialStringTok{\textquotesingle{}r{-}\textquotesingle{}}\OperatorTok{,} \SpecialStringTok{\textquotesingle{}LineWidth\textquotesingle{}}\OperatorTok{,} \FloatTok{2}\NormalTok{)}\OperatorTok{;}
\VariableTok{xlabel}\NormalTok{(}\SpecialStringTok{\textquotesingle{}Voltage X\textquotesingle{}}\NormalTok{)}\OperatorTok{;} \VariableTok{ylabel}\NormalTok{(}\SpecialStringTok{\textquotesingle{}f\_X(x)\textquotesingle{}}\NormalTok{)}\OperatorTok{;}
\VariableTok{title}\NormalTok{(}\SpecialStringTok{\textquotesingle{}Input: Gaussian Noise Voltage\textquotesingle{}}\NormalTok{)}\OperatorTok{;}
\VariableTok{legend}\NormalTok{(}\SpecialStringTok{\textquotesingle{}Empirical\textquotesingle{}}\OperatorTok{,} \SpecialStringTok{\textquotesingle{}Theoretical\textquotesingle{}}\NormalTok{)}\OperatorTok{;}

\VariableTok{subplot}\NormalTok{(}\FloatTok{2}\OperatorTok{,}\FloatTok{1}\OperatorTok{,}\FloatTok{2}\NormalTok{)}\OperatorTok{;}
\VariableTok{histogram}\NormalTok{(}\VariableTok{Y}\OperatorTok{,} \FloatTok{100}\OperatorTok{,} \SpecialStringTok{\textquotesingle{}Normalization\textquotesingle{}}\OperatorTok{,} \SpecialStringTok{\textquotesingle{}pdf\textquotesingle{}}\NormalTok{)}\OperatorTok{;}
\VariableTok{hold} \VariableTok{on}\OperatorTok{;}
\CommentTok{\% Theoretical: Y \textasciitilde{} chi{-}squared with 1 degree of freedom}
\VariableTok{y\_theory} \OperatorTok{=} \VariableTok{linspace}\NormalTok{(}\FloatTok{0.01}\OperatorTok{,} \FloatTok{10}\OperatorTok{,} \FloatTok{500}\NormalTok{)}\OperatorTok{;}
\VariableTok{pdf\_y\_theory} \OperatorTok{=} \VariableTok{chi2pdf}\NormalTok{(}\VariableTok{y\_theory}\OperatorTok{,} \FloatTok{1}\NormalTok{)}\OperatorTok{;}
\VariableTok{plot}\NormalTok{(}\VariableTok{y\_theory}\OperatorTok{,} \VariableTok{pdf\_y\_theory}\OperatorTok{,} \SpecialStringTok{\textquotesingle{}r{-}\textquotesingle{}}\OperatorTok{,} \SpecialStringTok{\textquotesingle{}LineWidth\textquotesingle{}}\OperatorTok{,} \FloatTok{2}\NormalTok{)}\OperatorTok{;}
\VariableTok{xlabel}\NormalTok{(}\SpecialStringTok{\textquotesingle{}Power Y = X\^{}2\textquotesingle{}}\NormalTok{)}\OperatorTok{;} \VariableTok{ylabel}\NormalTok{(}\SpecialStringTok{\textquotesingle{}f\_Y(y)\textquotesingle{}}\NormalTok{)}\OperatorTok{;}
\VariableTok{title}\NormalTok{(}\SpecialStringTok{\textquotesingle{}Output: Instantaneous Power\textquotesingle{}}\NormalTok{)}\OperatorTok{;}
\VariableTok{legend}\NormalTok{(}\SpecialStringTok{\textquotesingle{}Empirical\textquotesingle{}}\OperatorTok{,} \SpecialStringTok{\textquotesingle{}Theoretical (χ²₁)\textquotesingle{}}\NormalTok{)}\OperatorTok{;}
\end{Highlighting}
\end{Shaded}

\PSsubsection{Example 5: PMF vs PDF vs CDF Comparison}

\tcbset{PScodemode/.style={breakable}}

\begin{Shaded}
\begin{Highlighting}[]
\CommentTok{\%\% Complete Visualization: PMF/PDF and CDF side by side}
\VariableTok{figure}\OperatorTok{;}

\CommentTok{\% {-}{-}{-} Discrete Case: Number of retransmissions {-}{-}{-}}
\CommentTok{\% Geometric distribution: number of trials until first success}
\VariableTok{p\_success} \OperatorTok{=} \FloatTok{0.3}\OperatorTok{;}   \CommentTok{\% success probability per trial}
\VariableTok{k} \OperatorTok{=} \FloatTok{1}\OperatorTok{:}\FloatTok{15}\OperatorTok{;}

\CommentTok{\% PMF}
\VariableTok{subplot}\NormalTok{(}\FloatTok{2}\OperatorTok{,}\FloatTok{2}\OperatorTok{,}\FloatTok{1}\NormalTok{)}\OperatorTok{;}
\VariableTok{pmf\_geom} \OperatorTok{=} \VariableTok{geopdf}\NormalTok{(}\VariableTok{k}\OperatorTok{{-}}\FloatTok{1}\OperatorTok{,} \VariableTok{p\_success}\NormalTok{)}\OperatorTok{;}  \CommentTok{\% MATLAB\textquotesingle{}s geopdf counts failures}
\VariableTok{stem}\NormalTok{(}\VariableTok{k}\OperatorTok{,} \VariableTok{pmf\_geom}\OperatorTok{,} \SpecialStringTok{\textquotesingle{}filled\textquotesingle{}}\OperatorTok{,} \SpecialStringTok{\textquotesingle{}LineWidth\textquotesingle{}}\OperatorTok{,} \FloatTok{1.5}\NormalTok{)}\OperatorTok{;}
\VariableTok{xlabel}\NormalTok{(}\SpecialStringTok{\textquotesingle{}k (transmissions)\textquotesingle{}}\NormalTok{)}\OperatorTok{;} \VariableTok{ylabel}\NormalTok{(}\SpecialStringTok{\textquotesingle{}P(X = k)\textquotesingle{}}\NormalTok{)}\OperatorTok{;}
\VariableTok{title}\NormalTok{(}\SpecialStringTok{\textquotesingle{}Discrete PMF: Retransmissions\textquotesingle{}}\NormalTok{)}\OperatorTok{;}
\VariableTok{grid} \VariableTok{on}\OperatorTok{;}

\CommentTok{\% CDF}
\VariableTok{subplot}\NormalTok{(}\FloatTok{2}\OperatorTok{,}\FloatTok{2}\OperatorTok{,}\FloatTok{2}\NormalTok{)}\OperatorTok{;}
\VariableTok{cdf\_geom} \OperatorTok{=} \VariableTok{geocdf}\NormalTok{(}\VariableTok{k}\OperatorTok{{-}}\FloatTok{1}\OperatorTok{,} \VariableTok{p\_success}\NormalTok{)}\OperatorTok{;}
\VariableTok{stairs}\NormalTok{(}\VariableTok{k}\OperatorTok{,} \VariableTok{cdf\_geom}\OperatorTok{,} \SpecialStringTok{\textquotesingle{}LineWidth\textquotesingle{}}\OperatorTok{,} \FloatTok{2}\NormalTok{)}\OperatorTok{;}
\VariableTok{xlabel}\NormalTok{(}\SpecialStringTok{\textquotesingle{}k\textquotesingle{}}\NormalTok{)}\OperatorTok{;} \VariableTok{ylabel}\NormalTok{(}\SpecialStringTok{\textquotesingle{}P(X ≤ k)\textquotesingle{}}\NormalTok{)}\OperatorTok{;}
\VariableTok{title}\NormalTok{(}\SpecialStringTok{\textquotesingle{}Discrete CDF\textquotesingle{}}\NormalTok{)}\OperatorTok{;}
\VariableTok{grid} \VariableTok{on}\OperatorTok{;} \VariableTok{ylim}\NormalTok{([}\FloatTok{0} \FloatTok{1.1}\NormalTok{])}\OperatorTok{;}

\CommentTok{\% {-}{-}{-} Continuous Case: Delay time {-}{-}{-}}
\CommentTok{\% Exponential distribution}
\VariableTok{mu} \OperatorTok{=} \FloatTok{2}\OperatorTok{;}  \CommentTok{\% mean delay in ms}
\VariableTok{t} \OperatorTok{=} \VariableTok{linspace}\NormalTok{(}\FloatTok{0}\OperatorTok{,} \FloatTok{10}\OperatorTok{,} \FloatTok{500}\NormalTok{)}\OperatorTok{;}

\CommentTok{\% PDF}
\VariableTok{subplot}\NormalTok{(}\FloatTok{2}\OperatorTok{,}\FloatTok{2}\OperatorTok{,}\FloatTok{3}\NormalTok{)}\OperatorTok{;}
\VariableTok{pdf\_exp} \OperatorTok{=} \VariableTok{exppdf}\NormalTok{(}\VariableTok{t}\OperatorTok{,} \VariableTok{mu}\NormalTok{)}\OperatorTok{;}
\VariableTok{plot}\NormalTok{(}\VariableTok{t}\OperatorTok{,} \VariableTok{pdf\_exp}\OperatorTok{,} \SpecialStringTok{\textquotesingle{}LineWidth\textquotesingle{}}\OperatorTok{,} \FloatTok{2}\NormalTok{)}\OperatorTok{;}
\VariableTok{xlabel}\NormalTok{(}\SpecialStringTok{\textquotesingle{}Delay (ms)\textquotesingle{}}\NormalTok{)}\OperatorTok{;} \VariableTok{ylabel}\NormalTok{(}\SpecialStringTok{\textquotesingle{}f\_X(x) [ms\^{}\{{-}1\}]\textquotesingle{}}\NormalTok{)}\OperatorTok{;}
\VariableTok{title}\NormalTok{(}\SpecialStringTok{\textquotesingle{}Continuous PDF: Delay\textquotesingle{}}\NormalTok{)}\OperatorTok{;}
\VariableTok{grid} \VariableTok{on}\OperatorTok{;}

\CommentTok{\% CDF}
\VariableTok{subplot}\NormalTok{(}\FloatTok{2}\OperatorTok{,}\FloatTok{2}\OperatorTok{,}\FloatTok{4}\NormalTok{)}\OperatorTok{;}
\VariableTok{cdf\_exp} \OperatorTok{=} \VariableTok{expcdf}\NormalTok{(}\VariableTok{t}\OperatorTok{,} \VariableTok{mu}\NormalTok{)}\OperatorTok{;}
\VariableTok{plot}\NormalTok{(}\VariableTok{t}\OperatorTok{,} \VariableTok{cdf\_exp}\OperatorTok{,} \SpecialStringTok{\textquotesingle{}LineWidth\textquotesingle{}}\OperatorTok{,} \FloatTok{2}\NormalTok{)}\OperatorTok{;}
\VariableTok{xlabel}\NormalTok{(}\SpecialStringTok{\textquotesingle{}Delay (ms)\textquotesingle{}}\NormalTok{)}\OperatorTok{;} \VariableTok{ylabel}\NormalTok{(}\SpecialStringTok{\textquotesingle{}P(X ≤ x)\textquotesingle{}}\NormalTok{)}\OperatorTok{;}
\VariableTok{title}\NormalTok{(}\SpecialStringTok{\textquotesingle{}Continuous CDF\textquotesingle{}}\NormalTok{)}\OperatorTok{;}
\VariableTok{grid} \VariableTok{on}\OperatorTok{;} \VariableTok{ylim}\NormalTok{([}\FloatTok{0} \FloatTok{1.1}\NormalTok{])}\OperatorTok{;}

\VariableTok{sgtitle}\NormalTok{(}\SpecialStringTok{\textquotesingle{}PMF/PDF vs CDF: Discrete and Continuous\textquotesingle{}}\OperatorTok{,} \SpecialStringTok{\textquotesingle{}FontSize\textquotesingle{}}\OperatorTok{,} \FloatTok{14}\NormalTok{)}\OperatorTok{;}
\end{Highlighting}
\end{Shaded}

\PSsection{2.11}{Practice Problems}

\begin{PSbox}[unbreakable]{prob}{Problem 1: Identifying RV Type}

Classify each as discrete or continuous and state the range:

\begin{enumerate}
\def\labelenumi{(\alph{enumi})}
\tightlist
\item
  Number of dropped calls in one hour
\item
  Received signal strength (dBm)
\item
  Number of bits in error in a 1500-byte Ethernet frame
\item
  Time between packet arrivals
\item
  Quantizer output level (3-bit quantizer)
\end{enumerate}

\PSsolution{} 

\begin{enumerate}
\def\labelenumi{(\alph{enumi})}
\tightlist
\item
  Discrete, $X \in \{0,\, 1,\, 2,\, \ldots \}$
\item
  Continuous, $X \in (-\infty,\, \infty)$ practically $\sim (-120,\, 0)$ dBm
\item
  Discrete, $X \in \{0,\, 1,\, 2,\, \ldots,\, 12000\}$
\item
  Continuous, $T \in [0,\, \infty)$
\item
  Discrete, $X \in \{0,\, 1,\, 2,\, 3,\, 4,\, 5,\, 6,\, 7\}$
\end{enumerate}

\end{PSbox}

\begin{PSbox}[unbreakable]{prob}{Problem 2: PMF Properties}

A random variable $X$ has PMF: $p_{X}(0) = 0.1,\, p_{X}(1) = 0.3,\, p_{X}(2) = c,\, p_{X}(3) = 0.2$.

\begin{enumerate}
\def\labelenumi{(\alph{enumi})}
\tightlist
\item
  Find $c$.
\item
  Find $P(X \ge 2)$.
\item
  Find $F_{X}(1.5)$.
\item
  Find $P(0.5 < X \le 2.5)$.
\end{enumerate}

\PSsolution{} 

\begin{enumerate}
\def\labelenumi{(\alph{enumi})}
\item
  Sum = 1: $0.1 + 0.3 + c + 0.2 = 1 \to c = 0.4$
\item
  $P(X \ge 2) = p_{X}(2) + p_{X}(3) = 0.4 + 0.2 = 0.6$
\item
  $F_{X}(1.5) = P(X \le 1.5) = p_{X}(0) + p_{X}(1) = 0.1 + 0.3 = 0.4$
\item
  $P(0.5 < X \le 2.5) = p_{X}(1) + p_{X}(2) = 0.3 + 0.4 = 0.7$
\end{enumerate}

\end{PSbox}

\begin{PSbox}[unbreakable]{prob}{Problem 3: PDF and CDF}

A random variable has PDF:\PSbr{}$f_{X}(x)$ = cx for $0 \le x \le 4$, and 0 otherwise.

\begin{enumerate}
\def\labelenumi{(\alph{enumi})}
\tightlist
\item
  Find $c$.
\item
  Find the CDF $F_{X}(x)$.
\item
  Find $P(1 \le X \le 3)$.
\item
  Find the value $x_{0}$ such that $P(X \le x_{0}) = 0.5$.
\end{enumerate}

\PSsolution{} 

\begin{enumerate}
\def\labelenumi{(\alph{enumi})}
\item
  $\int_{0}^{4}$ cx $dx = c \cdot \left[\frac{x^{2}}{2}\right]_{0}^{4} = c \cdot 8 = 1 \to c = \frac{1}{8}$
\item
  $F_{X}(x) = \int_{0}^{x} \frac{t}{8}dt = \frac{x^{2}}{16}$, for $0 \le x \le 4$\PSbr{}$F_{X}(x) = 0$ for $x < 0$; $F_{X}(x) = 1$ for $x > 4$
\item
  $\displaystyle P(1 \le X \le 3) = F_{X}(3) - F_{X}(1) = \frac{9}{16} - \frac{1}{16} = \frac{8}{16} = 0.5$
\item
  $\displaystyle \frac{x_{0}^{2}}{16} = 0.5 \to x_{0}^{2} = 8 \to x_{0} = 2\sqrt{2} \approx 2.83$
\end{enumerate}

\end{PSbox}

\begin{PSbox}[unbreakable]{prob}{Problem 4: Transformation}

A sensor measures temperature $X$ uniformly distributed on [20, 30] \ensuremath{^{\circ}}C.\PSbr{}The output voltage is $Y = 0.1X - 2$ (linear sensor).

\begin{enumerate}
\def\labelenumi{(\alph{enumi})}
\tightlist
\item
  Find the PDF of $Y$.
\item
  Find the range of $Y$.
\item
  Find $P(Y > 0.5)$.
\end{enumerate}

\PSsolution{} 

\begin{enumerate}
\def\labelenumi{(\alph{enumi})}
\item
  $X \sim \mathrm{Uniform}[20,\, 30]$, so $f_{X}(x) = \frac{1}{10}$ for $20 \le x \le 30$.\PSbr{}$Y = 0.1X - 2$, so $a = 0.1,\, b = -2$.\PSbr{}$f_{Y}(y) = \frac{1}{\lvert a\rvert} \cdot f_{X}\left(\frac{y-b}{a}\right) = \frac{1}{0.1} \cdot \frac{1}{10} = 1$ for the valid range.
\item
  When $X = 20$: $Y = 0.1(20) - 2 = 0$; When $X = 30$: $Y = 0.1(30) - 2 = 1$\PSbr{}Range of $Y$: [0, 1]. So $Y \sim \mathrm{Uniform}[0,\, 1]$.
\item
  $P(Y > 0.5) = 1 - F_{Y}(0.5) = 1 - 0.5 = 0.5$
\end{enumerate}

\end{PSbox}

\begin{PSbox}[unbreakable]{prob}{Problem 5: Non-Monotonic Transformation}

A voltage $X \sim \mathrm{Uniform}[-1,\, 1] V$. The instantaneous power is $Y = X^{2}$.

\begin{enumerate}
\def\labelenumi{(\alph{enumi})}
\tightlist
\item
  Find the CDF of $Y$.
\item
  Find the PDF of $Y$.
\item
  Verify using MATLAB simulation.
\end{enumerate}

\PSsolution{} 

\begin{enumerate}
\def\labelenumi{(\alph{enumi})}
\tightlist
\item
  For $0 \le y \le 1$:\PSbr{}$F_{Y}(y) = P(X^{2} \le y) = P\left(-\sqrt{y} \le X \le \sqrt{y}\right) = F_{X}\left(\sqrt{y}\right) - F_{X}\left(-\sqrt{y}\right)$
\end{enumerate}

Since $X \sim \mathrm{Uniform}[-1,\,1]$: $F_{X}(x) = \frac{x+1}{2}$

$\displaystyle F_{Y}(y) = \frac{\sqrt{y} + 1}{2} - \frac{-\sqrt{y} + 1}{2} = \sqrt{y}$

\begin{enumerate}
\def\labelenumi{(\alph{enumi})}
\setcounter{enumi}{1}
\item
  $f_{Y}(y) = \frac{dF_{Y}}{dy} = \frac{1}{2\sqrt{y}}$ for $0 < y \le 1$
\item
  MATLAB:
\end{enumerate}

\tcbset{PScodemode/.style={unbreakable}}

\begin{Shaded}
\begin{Highlighting}[]
\VariableTok{N} \OperatorTok{=} \FloatTok{100000}\OperatorTok{;}
\VariableTok{X} \OperatorTok{=} \FloatTok{2}\OperatorTok{*}\VariableTok{rand}\NormalTok{(}\FloatTok{1}\OperatorTok{,}\VariableTok{N}\NormalTok{) }\OperatorTok{{-}} \FloatTok{1}\OperatorTok{;}   \CommentTok{\% Uniform[{-}1,1]}
\VariableTok{Y} \OperatorTok{=} \VariableTok{X}\OperatorTok{.\^{}}\FloatTok{2}\OperatorTok{;}
\VariableTok{histogram}\NormalTok{(}\VariableTok{Y}\OperatorTok{,} \FloatTok{100}\OperatorTok{,} \SpecialStringTok{\textquotesingle{}Normalization\textquotesingle{}}\OperatorTok{,} \SpecialStringTok{\textquotesingle{}pdf\textquotesingle{}}\NormalTok{)}\OperatorTok{;}
\VariableTok{hold} \VariableTok{on}\OperatorTok{;}
\VariableTok{y} \OperatorTok{=} \VariableTok{linspace}\NormalTok{(}\FloatTok{0.01}\OperatorTok{,} \FloatTok{1}\OperatorTok{,} \FloatTok{200}\NormalTok{)}\OperatorTok{;}
\VariableTok{plot}\NormalTok{(}\VariableTok{y}\OperatorTok{,} \FloatTok{1}\OperatorTok{./}\NormalTok{(}\FloatTok{2}\OperatorTok{*}\VariableTok{sqrt}\NormalTok{(}\VariableTok{y}\NormalTok{))}\OperatorTok{,} \SpecialStringTok{\textquotesingle{}r{-}\textquotesingle{}}\OperatorTok{,} \SpecialStringTok{\textquotesingle{}LineWidth\textquotesingle{}}\OperatorTok{,} \FloatTok{2}\NormalTok{)}\OperatorTok{;}
\VariableTok{legend}\NormalTok{(}\SpecialStringTok{\textquotesingle{}Empirical\textquotesingle{}}\OperatorTok{,} \SpecialStringTok{\textquotesingle{}Theory: 1/(2√y)\textquotesingle{}}\NormalTok{)}\OperatorTok{;}
\VariableTok{xlabel}\NormalTok{(}\SpecialStringTok{\textquotesingle{}Y = X²\textquotesingle{}}\NormalTok{)}\OperatorTok{;} \VariableTok{ylabel}\NormalTok{(}\SpecialStringTok{\textquotesingle{}f\_Y(y)\textquotesingle{}}\NormalTok{)}\OperatorTok{;}
\end{Highlighting}
\end{Shaded}

\end{PSbox}

\PSsection{2.12}{Key Takeaways}

\begin{enumerate}
\def\labelenumi{\arabic{enumi}.}
\tightlist
\item
  \textbf{Random variables} convert experimental outcomes into numbers for mathematical analysis
\item
  \textbf{PMF} gives the actual probability for discrete RVs (bar heights sum to 1)
\item
  \textbf{PDF} gives probability density for continuous RVs (area gives probability, not height)
\item
  \textbf{CDF} works universally — gives cumulative probability $P(X \le x)$
\item
  \textbf{Transformations} let us derive distributions of functions of random variables
\item
  The \textbf{CDF method} (find CDF first, then differentiate) is often the safest approach for transformations
\end{enumerate}

\emph{Next Module: {Module 3 — Important Probability Distributions}}
